\documentclass[10pt]{article}

\usepackage{amsmath,amssymb,amsthm}

\newtheorem{definition}{Definition}
\newtheorem{problem}{Problem}
\newtheorem{theorem}{Theorem}
\newtheorem{lemma}[theorem]{Lemma}
\newtheorem{corollary}[theorem]{Corollary}
\newtheorem{remark}[theorem]{Remark}

\newcommand{\SKF}{\text{SK1Flow}}
\newcommand{\arc}[2]{#1\!\to\!#2}
\DeclareMathOperator{\capa}{cap}

\title{A Sub-Cubic Bound for the SK1Flow Maximum-Flow Algorithm - Medium}
\author{Nikita Gaevoy}
\date{}

\begin{document}

\maketitle

\section{Preliminaries}

\subsection{Excess networks and preflows}

An \emph{excess network} $N=(V,A,c,e_0,t)$ consists of a vertex set $V$ with
$|V|=n$; an arc set $A$, partitioned into $m$ \emph{pairs} $\{a,a^{R}\}$ with
$a=\arc{v}{u}$ and $a^{R}=\arc{u}{v}$; a non-negative capacity
$c\colon A\to T_{\ge0}$; an \emph{initial excess} $e_0\colon V\to T_{\ge0}$; and
a \emph{sink} $t\in V$. Arcs may repeat, so the underlying graph is a directed
multigraph. Here $T$ is a totally ordered abelian group whose order is
translation invariant,
\[
x\le y\implies x+z\le y+z,
\]
for instance $\mathbb{Z}$, $\mathbb{Q}$, or arbitrary-precision integers. A
one-directional arc of capacity $\gamma$ is the pair $(\gamma,0)$, and $m$ is
the number of pairs, so $\sum_{v}\deg v=2m$ where $\deg v$ counts the arcs
leaving $v$ in either orientation.

There is no source vertex: sources are the vertices carrying positive initial
excess, which makes multiple sources free of charge. Throughout, $r(a)$ denotes
the current \emph{residual capacity} of the arc $a$ and $e(v)$ the current
excess of $v$; both start at $c$ and $e_0$ and are modified only by the
operation
\[
  \mathrm{push}(\arc{v}{u})\colon\quad
  \delta:=\min\bigl(e(v),r(\arc{v}{u})\bigr),\qquad
  \begin{aligned}
    r(\arc{v}{u})&\mathrel{-}=\delta, & r(\arc{u}{v})&\mathrel{+}=\delta,\\
    e(v)&\mathrel{-}=\delta,          & e(u)&\mathrel{+}=\delta.
  \end{aligned}
\]
The push is \emph{positive} if $\delta>0$. Two invariants are immediate:
$r(a)+r(a^{R})$ is constant, and $r\ge0$ and $e\ge0$ hold at all times. The
intermediate states are therefore \emph{preflows}, not flows: excess may
accumulate at intermediate vertices and stay there.

\begin{definition}[Residual graph, distance]\label{def:residual}
The \emph{residual graph} $G_r$ has vertex set $V$ and contains the arc
$\arc{v}{u}$ exactly when $r(\arc{v}{u})>0$. We write $d(v)$ for the number of
arcs on a shortest $v\rightsquigarrow t$ path in $G_r$, and $d(v)=\infty$ when
$t$ is unreachable from $v$ in $G_r$.
\end{definition}

The quantity to be computed is the maximum amount deliverable to $t$: the
maximum, over all feasible routings of the initial excess, of the amount that
ends at $t$. Equivalently it is the value of a maximum flow from a super-source
$s$ joined to each $v$ by an arc of capacity $e_0(v)$.

\subsection{The algorithm}

$\SKF$ is a phase algorithm. Each phase computes the exact distances $d$ by one
backward breadth-first search from the sink, and then makes a single greedy
sweep over the vertices in order of decreasing distance, each vertex forwarding
whatever excess it holds one layer down.

\begin{definition}[Heights, admissibility]\label{def:height}
The \emph{height} of $v$ in a phase is $\ell(v)=d(v)+1$ if $d(v)<\infty$, and
$\ell(v)=\infty$ otherwise, where $d$ is taken at the start of the phase. The
offset of $1$ is a convenience: it puts the sink at height $1$ and leaves
height $0$ empty, so the sink has no admissible out-arc and never pushes flow
back out. An arc $\arc{v}{u}$ is \emph{admissible in the phase} if
$\ell(u)=\ell(v)-1$ with both heights finite and $r(\arc{v}{u})>0$.
\end{definition}

\begin{definition}[$\SKF$]\label{def:alg}
Given an excess network $N$ and its sink $t$, repeat the following
\emph{phase}.
\begin{enumerate}
  \item \emph{Layering.} Run a breadth-first search from $t$ in the reverse of
  $G_r$. This sets $\ell(v)=d(v)+1$ for every $v$ with $d(v)<\infty$ and
  $\ell(v)=\infty$ for the rest, and produces a list
  $\mathrm{order}$ of exactly the vertices of finite height, without
  repetition, in non-decreasing order of $\ell$.

  \item \emph{Sweep.} Traverse $\mathrm{order}$ backwards, that is in
  non-increasing order of $\ell$, so that every vertex of finite height is
  visited exactly once. On visiting $v$, scan the arcs leaving $v$ in a fixed
  but arbitrary order, stopping as soon as $e(v)=0$, and perform
  $\mathrm{push}(\arc{v}{u})$ on every admissible arc $\arc{v}{u}$
  encountered.

  \item \emph{Termination test.} If the sweep performed no positive push, stop.
\end{enumerate}
On stopping, return $e(t)$ and reset $e(t)$ to zero, so that a further call
resumes from the current residual state instead of restarting.
\end{definition}

Two features distinguish $\SKF$ from its classical relatives, and both are
load-bearing. Its intermediate states are preflows, so a sweep may strand
excess at an intermediate vertex, whereas a blocking-flow phase of Dinitz's
algorithm~\cite{Dinitz70} returns the surplus and ends with a genuine flow.
At the same time it recomputes \emph{exact} distances once per phase, whereas
push-relabel~\cite{GoldbergTarjan88} maintains lazy lower-bound estimates and
relabels a single vertex at a time. $\SKF$ is the remaining combination:
preflow-shaped, with exact per-phase relabelling.
Removing the stranding --- a backward pass returning each vertex's surplus along
the arcs that delivered it --- is exactly the balancing performed inside one
layered network by the preflow-based blocking-flow routines of
Karzanov~\cite{Karzanov74} and of Malhotra, Kumar and
Maheshwari~\cite{MKM78}; with the exact per-phase relabelling already in place a
phase then becomes a full Dinitz phase, and the algorithm becomes Dinitz's. The
excess carried across phases is not an incidental feature of the implementation,
and it is what the problem below is about.

\begin{remark}\label{rem:order}
The arc order inside a sweep is arbitrary and is not part of the problem: all
the bounds below hold for every choice. It does change \emph{which} vertices
strand excess, hence the constants, so a solution may not assume a particular
order unless it establishes the bound for all of them.
\end{remark}

\begin{remark}\label{rem:correct}
$\SKF$ is correct. On termination let $\mathcal{R}=\{v:d(v)<\infty\}$. If a
phase makes no positive push then no $v\ne t$ with $e(v)>0$ lies in
$\mathcal{R}$, since such a $v$ would have an admissible first arc on a shortest
$v\rightsquigarrow t$ path and would have pushed along it. Hence every arc
entering $\mathcal{R}$ from outside is saturated, and summing the excess
conservation identity over $\mathcal{R}$ gives
$e(t)=\sum_{v\in\mathcal{R}}e_0(v)+\capa(V\setminus\mathcal{R}\to\mathcal{R})$,
which is an upper bound on
what any routing can deliver. The excess stranded outside $\mathcal{R}$ could not have
reached $t$ under any routing, so no final balancing stage is needed; the
residual state left behind is a preflow and must not be read as a flow
decomposition without first cancelling that excess.
\end{remark}

\subsection{The cost measure}

Cost is measured in elementary operations: one comparison, one addition, one
subtraction, or one copy of an element of $T$, and one word-RAM operation on an
index or a height. For $T=\mathbb{Z}$ inside a machine word these are $O(1)$
each. Every bound below is strongly polynomial, that is, independent of the
capacity values.

\begin{definition}[The three counters]\label{def:counters}
Over a whole run of $\SKF$ on an excess network $N$ write
\begin{itemize}
  \item $P=P(N)$ for the number of phases, including the final pushless one;
  \item $S=S(N)$ for the number of positive pushes that leave their arc
  saturated, $r(\arc{v}{u})=0$;
  \item $Q=Q(N)$ for the number of positive pushes that leave their tail empty,
  $e(v)=0$.
\end{itemize}
A push that does both is counted in $S$ and in $Q$. Every positive push lies in
$S\cup Q$, since $\delta=\min(e(v),r(\arc{v}{u}))$ attains one of its
arguments.
\end{definition}

The counters $P$, $S$, $Q$ depend only on $N$ and on the arc order of
Remark~\ref{rem:order}; they are properties of the run, not of any particular
implementation of step~1 of Definition~\ref{def:alg}.

\begin{remark}[The layering step is not the subject]\label{rem:layering}
Step~1 as written costs $\Theta(n+m)$ per phase whatever the phase does, and a
literal implementation therefore runs in $\Theta(P(n+m))$. This is an artefact
of recomputing from scratch, and the problem below is posed for the
\emph{algorithm}, not for that implementation: a solver may replace step~1 by
any procedure that produces the same heights $\ell=d+1$, and in particular by
an incremental relabeller. Theorem~\ref{thm:time} records what this buys,
and Problem~\ref{prob:upper} restates the licence explicitly. The sweep, the
pushes, and hence the counters of Definition~\ref{def:counters} are unaffected
by any such substitution.
\end{remark}

\begin{theorem}[Running time]\label{thm:time}
Maintain the heights incrementally, recomputing only the vertices whose exact
distance changes, and resume each vertex's arc scan from a current-arc pointer
that is reset only when its height changes. Then $\SKF$ runs in
\[
  O(n+m+nm+S+Q)\;=\;O(n+nm+Q)
\]
elementary operations, and no implementation runs in fewer than
$\Omega(n+m+S+Q)$. Consequently $\SKF$ runs in $O(n^{3})$ unconditionally, and,
on networks with $m\ge1$, in $O(nm)$ if and only if $Q=O(nm)$.
\end{theorem}

\begin{proof}
A finite height is at most $n$ and never decreases, so each vertex's height
changes at most $n$ times; each change costs $O(\deg v)$ to propagate and is the
only event that resets that vertex's pointer, so the relabelling and the arc
scans together cost $O\bigl(\sum_v n\deg v\bigr)=O(nm)$. An arc that has been
saturated regains residual capacity only through a push on its partner, which
raises the height of its head by at least two, so each arc is saturated $O(n)$
times and $S=O(nm)$. Every positive push is counted by $S+Q$ and costs $O(1)$,
and visiting the active vertices in non-increasing height order costs $O(1)$
amortised per visit with a monotone bucket queue over the $n+1$ possible
heights; the initial arrays and search cost $O(n+m)$. For the lower bound, every
implementation must read its input and perform the pushes of step~2, one
elementary operation each.

For the unconditional bound, each vertex empties its tail at most once per phase,
so $Q\le nP$; and every phase but the last two strictly raises $d(v)$ for some
$v$, while $\sum_v d(v)$ can rise at most $n(n-1)$ times, so $P=O(n^{2})$ and
$Q=O(n^{3})$. Conversely an $O(nm)$ running time forces $Q=O(nm)$, since $Q$ is
a lower bound on the number of operations.
\end{proof}

\section{Problem formulation}

\begin{problem}[Upper bound]\label{prob:upper}
Prove that $\SKF$ runs in $o(n^{3})$ elementary operations. Formally: exhibit a
function $T(n,m)$ with
\[
  T(n,m)=o(n^{3})\quad\text{as }n\to\infty\text{ for some infinite family of
  densities }m=m(n),
\]
and an implementation of $\SKF$ that runs in $O\bigl(T(n,m)\bigr)$ elementary
operations on every excess network with $n$ vertices and $m$ arc pairs.

The implementation is free to compute the heights of step~1 of
Definition~\ref{def:alg} by any means, as long as the heights it produces are
the exact ones, $\ell=d+1$; the sweep and the pushes of step~2 must be
performed exactly as stated.

On densities $m=O(n^{2-\varepsilon})$, Theorem~\ref{thm:time} makes this
equivalent to the purely combinatorial statement
\[
  Q=o(n^{3})
\]
about the run: every other term of the running time is already $o(n^{3})$
there, while conversely $Q$ is a lower bound on the operations of any
implementation.
\end{problem}

\begin{remark}[Why $o(n^{3})$ rather than $O(nm)$]\label{rem:density}
On dense graphs the two coincide: $nm=\Theta(n^{3})$ when $m=\Theta(n^{2})$, so
the conjectured $O(nm)$ is \emph{not} an improvement there. Nor can the
accounting of Theorem~\ref{thm:time} deliver one, whatever is proved about $Q$:
its other terms already amount to $\Theta(nm)$, and beating $O(n^{3})$ on dense
graphs would need a cheaper way to maintain the heights than the $O(nm)$ of
Theorem~\ref{thm:time}, which nothing here rules out but which is a separate
question. Problem~\ref{prob:upper} is therefore stated as $o(n^{3})$
for \emph{some} infinite family of densities, and the interesting regime is
$m=o(n^{2})$. Three targets of increasing strength are worth naming:
\begin{enumerate}
  \item $Q=o(n^{3})$ for sparse graphs, $m=O(n)$, where the current bound
  $Q=O(n^{3})$ is a factor $n$ away from the conjecture;
  \item $Q=O(n^{2}\sqrt m\,)$, the bound that holds for the genuine
  highest-label push-relabel rule~\cite{CheriyanMaheshwari89} and which is
  $o(n^{3})$ whenever $m=o(n^{2})$. The sweep of step~2 visits vertices in
  non-increasing height order, which makes $\SKF$ resemble that rule; it is not
  that rule, because it defers every relabel to the end of the phase instead of
  lifting a stuck vertex at once, and whether the potential argument
  of~\cite{CheriyanMaheshwari89,CheriyanMehlhorn99} survives the deferral is
  open;
  \item $Q=O(nm)$, the conjecture, which by Theorem~\ref{thm:time} gives
  an $O(nm)$ running time outright.
\end{enumerate}
Any one of them solves Problem~\ref{prob:upper}.
\end{remark}

\begin{remark}\label{rem:dichotomy}
A separate problem asks for a bound from the
other side: an infinite family of excess networks on which $\SKF$ performs
$\omega(nm)$ tail-emptying pushes, and hence takes $\omega(nm)$ time whatever
implementation of step~1 is used. The two are not alternatives and not a
dichotomy. The gap they straddle runs from $nm$ to $n^{3}$, and most bounds
inside it satisfy both at once: a proof that $Q=\Theta(n^{2}\sqrt m\,)$ on sparse
graphs, say, would solve this problem and the companion together. Only the
strongest form named in Remark~\ref{rem:density}, $Q=O(nm)$, excludes the
companion.
\end{remark}

\begin{remark}[Scope]\label{rem:scope}
The problem is about $\SKF$ as defined, not about the maximum flow problem.
Maximum flow is solvable in $O(nm)$ time~\cite{Orlin13} and in
$m^{1+o(1)}$ time~\cite{CKLPGS22}, so it is not asking for a new record; what is
at stake is the analysis of one short algorithm.

A solution must hold for every arc order of Remark~\ref{rem:order}, since the
implementation fixes none.
\end{remark}

\newcommand{\etalchar}[1]{$^{#1}$}
\begin{thebibliography}{CKL{\etalchar{+}}22}

\bibitem[CKL{\etalchar{+}}22]{CKLPGS22}
Li~Chen, Rasmus Kyng, Yang~P. Liu, Richard Peng, Maximilian Probst~Gutenberg,
  and Sushant Sachdeva.
\newblock Maximum flow and minimum-cost flow in almost-linear time.
\newblock In {\em Proceedings of the 63rd IEEE Annual Symposium on Foundations
  of Computer Science (FOCS 2022)}, pages 612--623, 2022.

\bibitem[CM89]{CheriyanMaheshwari89}
Joseph Cheriyan and S.~N. Maheshwari.
\newblock Analysis of preflow push algorithms for maximum network flow.
\newblock {\em SIAM Journal on Computing}, 18(6):1057--1086, 1989.

\bibitem[CM99]{CheriyanMehlhorn99}
Joseph Cheriyan and Kurt Mehlhorn.
\newblock An analysis of the highest-level selection rule in the preflow-push
  max-flow algorithm.
\newblock {\em Information Processing Letters}, 69(5):239--242, 1999.

\bibitem[Din70]{Dinitz70}
E.~A. Dinitz.
\newblock Algorithm for solution of a problem of maximum flow in networks with
  power estimation.
\newblock {\em Soviet Mathematics Doklady}, 11:1277--1280, 1970.

\bibitem[GT88]{GoldbergTarjan88}
Andrew~V. Goldberg and Robert~E. Tarjan.
\newblock A new approach to the maximum-flow problem.
\newblock {\em Journal of the ACM}, 35(4):921--940, 1988.

\bibitem[Kar74]{Karzanov74}
A.~V. Karzanov.
\newblock Determining the maximal flow in a network by the method of preflows.
\newblock {\em Soviet Mathematics Doklady}, 15(2):434--437, 1974.

\bibitem[MKM78]{MKM78}
V.~M. Malhotra, M.~Pramodh Kumar, and S.~N. Maheshwari.
\newblock An {$O(|V|^3)$} algorithm for finding maximum flows in networks.
\newblock {\em Information Processing Letters}, 7(6):277--278, 1978.

\bibitem[Orl13]{Orlin13}
James~B. Orlin.
\newblock Max flows in {$O(nm)$} time, or better.
\newblock In {\em Proceedings of the 45th Annual ACM Symposium on Theory of
  Computing (STOC 2013)}, pages 765--774, 2013.

\end{thebibliography}

\end{document}
